\chapter*{Interludium: Die Lebensgeschichte der Sieben}
\addcontentsline{toc}{chapter}{Interludium: Die Lebensgeschichte der Sieben}
\markboth{INTERLUDIUM}{INTERLUDIUM}

\begin{oberflaeche}
Begleiten wir eine einzige Zahl durch alle ihre Wohnungen.

Die Sieben kommt in der Welt mit vierundzwanzig Plätzen zur Welt --
früher geht es nicht, die Welt davor endet bei fünf. Dort
heißt sie $013$. Es ist eine beengte Wohnung: Drei Ziffern, alle
tragen etwas bei, nichts ist Dekoration.

Dann zieht sie um, Welt für Welt, immer werttreu -- sie will ja
dieselbe Sieben bleiben. Und bei jedem Umzug muss sie sich neu
einrichten: In der Welt mit hundertzwanzig Plätzen heißt sie
$0012$, in der nächsten $00011$, dann $000010$. Wer die Umzüge
nebeneinanderlegt, sieht ein rastloses Muster -- die Ziffern wandern,
tauschen, schrumpfen. Kein Umzug ist wie der andere, und nichts an
der Schreibweise verrät auf den ersten Blick, dass es immer dieselbe
Zahl ist.

Bis zur Welt mit vierzigtausenddreihundertzwanzig Plätzen. Dort
passiert etwas Neues: Die Sieben passt zum ersten Mal vollständig in
die letzte Ziffer -- sie heißt jetzt $0000007$, und man sieht ihr
ihren Wert wieder an. Von diesem Umzug an ändert sich nie mehr
etwas. Jede weitere Welt hängt nur noch eine stille Null vorne an:
$00000007$, $000000007$, immer so weiter. Die Sieben ist
angekommen; wir sagen: Sie hat ihren Eigenhorizont erreicht.

Jede Zahl hat diese Biographie: eine rastlose Jugend, in der jeder
größere Raum eine neue Identitätsverhandlung bedeutet, und ein
ruhiges Erwachsenenalter, ab dem mehr Raum einfach nur mehr Raum
ist. Große Zahlen haben eine lange Jugend -- die Zahl
siebenhundertzwanzig zum Beispiel wird bis zur Welt Nummer
siebenhundertzwanzig
umgebaut, und auf diesem langen Weg trägt sie zeitweise Gestalten,
die ihr niemand zutraut: In der neunten Welt ist sie schlicht eine
einzelne Eins an der vierten Stelle von rechts, denn zehn mal neun
mal acht ist siebenhundertzwanzig. Genau diese Doppelnatur -- dem
Wert nach sperrig, der Gestalt nach schlicht -- hat sich später als
Schlüssel zu einem der Hauptsätze erwiesen.

Merken muss man sich aus alledem nur eines: Im Horiraum ist
„dieselbe Zahl“ keine Selbstverständlichkeit, sondern eine
Entscheidung. Man kann den Wert festhalten und die Gestalt opfern,
oder umgekehrt. Wie viel diese Entscheidung kostet, ist keine
Geschmacksfrage mehr, sondern ein Satz -- aber das gehört ins
nächste Kapitel.
\end{oberflaeche}
